\section*{第 \ref{ch:b2:diffraction} 章——夫琅禾费衍射与空间滤波}

\begin{solution}{exo:b2:diffraction:1}
$2\lambda L/b = \qty{25}{mm}$；次极大位于 $1.43\lambda L/b = \qty{18}{mm}$ 及
$2.46\lambda L/b = \qty{31}{mm}$ 处；菲涅耳数 $b^2/\lambda L = 8 \times 10^{-3}$，
远场条件满足。加上透镜后，$2\lambda f/b = \qty{6.3}{mm}$。
\end{solution}

\begin{solution}{exo:b2:diffraction:2}
$1.22\lambda/D$：人眼 \qty{2.2e-4}{rad}（\qty{46}{''}）；\qty{10}{cm}：\qty{1.4}{''}；哈勃
望远镜 \qty{0.058}{''}；碟形天线 $1.22 \times 0.06/64 = \qty{1.1e-3}{rad} = \qty{3.9}{'}$；照相机：
$\qty{1.1e-4}{rad} \times \qty{50}{mm} = \qty{5.6}{\micro m}$ 为半径，即一个 \qty{11}{\micro m} 的圆斑——
相当于三个像素：在 $f/8$ 时，衍射已经限制了传感器。
\end{solution}

\begin{solution}{exo:b2:diffraction:3}
$\lambda f/a = \qty{2}{mm}$；包络 $2\lambda f/b = \qty{20}{mm}$；十条条纹；级次
$p = 5, 10, \dots$（$a/b = 5$）缺失。
\end{solution}

\begin{solution}{exo:b2:diffraction:4}
$\theta = 1.22\lambda/D = \qty{3.9e-4}{rad}$：光斑直径 \qty{290}{km}；若口径为 \qty{3.5}{m}，
则为 \qty{170}{m}。反射器自身的衍射，即 $\lambda/\qty{1}{m}$，使回波散布于
\qty{600}{m} 范围内：发出的 $10^{17}$ 个光子中，只有寥寥几个返回。
\end{solution}

\begin{solution}{exo:b2:diffraction:5}
(a) $\int_{-b/2}^{b/2}\eu^{\iu kx\sin\theta}\dd x = b\,\operatorname{sinc}(kb\sin\theta/2)$。(b)
由 $\dd(\sin x/x)/\dd x = 0$ 得 $\tan x = x$；$x = 4.49 = 1.43\pi$，$\operatorname{sinc}^2 =
0.047$。(c) $0.903$。(d) 相位变为 $kx(\sin\theta - \sin\theta_0)$，
图样绕 $\theta_0$ 呈相同的形状。
\end{solution}

\begin{solution}{exo:b2:diffraction:6}
$0.61\lambda/\mathrm{NA}$：\qty{220}{nm}；\qty{170}{nm}；\qty{110}{nm}。放大倍数超过约 $500\times$
后，目镜只会放大模糊（空放大）。电子在
\qty{4}{pm}、$\mathrm{NA} \approx 0.01$ 时：\qty{0.2}{nm}——恰好能分辨原子。
\end{solution}

\begin{solution}{exo:b2:diffraction:7}
(a) 积分可按 $x$ 与 $y$ 分解。(b) $2\lambda f/b = \qty{5}{mm}$。(c) 圆盘的
弦在远离中心处短于直径：其有效宽度较小，
图样则较宽。(d) 十字形的四条衍射
尖峰。
\end{solution}

\begin{solution}{exo:b2:diffraction:8}
(a) $\sin\theta = \pm0.275$（\qty{16}{\degree}），$\pm0.55$（\qty{33}{\degree}）。(b) $\sin\alpha \ge 0.275$。
(c) 仅有 $0$、$\pm1$ 级：得到的是周期正确的正弦像，而非
清晰的条纹。(d) $\sin\alpha \ge \lambda/2d = 0.14$：为分辨率的两倍。
\end{solution}

\begin{solution}{exo:b2:diffraction:9}
(a) $2.44\lambda f/D = \qty{15}{\micro m}$；针孔为 \qty{23}{\micro m}。(b) $\lambda/\qty{50}{\micro m} =
\qty{1.3e-2}{rad}$：偏离轴 \qty{0.25}{mm} 处被挡住。(c) $1\%$；针孔只保留
低空间频率成分。(d) 偏离几个微米时，针孔会
截去光斑本身的一部分。
\end{solution}

\begin{solution}{exo:b2:diffraction:10}
(a) 振幅与无遮挡波叠加，而在轴外无遮挡波为零：
此处 $|\underline s_{\text{缝}}| = |\underline s_{\text{条}}|$。(b) 条纹间距 $\lambda L/b =
\qty{16}{mm}$；测量间距即可得 $b$。(c) 细条的影子位于
明亮的未衍射光束之中。(d) 直径为 $d$ 的液滴衍射如同小孔
：环位于 $1.22\lambda/d$ 处，红色在外侧；由半径可得 $d$。
\end{solution}

\begin{solution}{exo:b2:diffraction:11}
(a) $|1 + \iu\varphi|^2 = 1 + \varphi^2$：一阶近似下均匀。(b) 去掉轴向光后，
$|\iu\varphi|^2 = \varphi^2$：二阶小量，较暗。(c) $|\iu + \iu\varphi|^2 =
1 + 2\varphi + \dots$：线性。(d) 轴上光斑比衍射光更亮，
将其衰减可使两振幅相当，从而提高对比度。
\end{solution}

\begin{solution}{exo:b2:diffraction:12}
(a) 半宽为 $\lambda/W$；除以 $p/a\cos\theta$：$\Delta\lambda = \lambda a\cos\theta/pW \approx
\lambda/pN$。(b) 准直光束的直径为 $Df_{\text{准直}}/f_{\text{望远}}$；
光栅若较小，则会浪费光并损失分辨本领。(c) 分辨率
总是孔径上相位改变一个波长所对应的角度，
即 $\lambda/(\text{尺寸})$。(d) 其像
宽度比光栅的衍射宽度更宽时占主导，
这是通常情形。
\end{solution}

\begin{solution}{pb:b2:diffraction:1}
\textbf{1.} $1.22\lambda/D = \qty{2.2e-4}{rad}$；$\times\qty{17}{mm} = \qty{3.8}{\micro m}$：相当于两个视锥间距——
视网膜的采样正处于光学极限。

\textbf{2.} \qty{56}{\micro m}（位于 \qty{25}{cm} 处），\qty{2.2}{mm}（位于 \qty{10}{m} 处）；\qty{1}{mm} 位于
\qty{50}{cm} 处时是极限的十倍：很容易分辨。

\textbf{3.} \qty{9.6e-5}{rad}；全孔径下透镜的像差，以及
较粗且合并的视杆细胞，抵消了增益。

\textbf{4.} $1.5/2.2 \times 10^{-4} = \qty{7}{km}$。

\textbf{5.} 衍射变差，达到 \qty{6.7e-4}{rad}，但离焦眼睛的弥散圆
随孔径缩小：景深更占优势。

\textbf{6.} \qty{6.7e-7}{rad}（\qty{0.14}{''}）；$\times\qty{10}{m} = \qty{6.7}{\micro m}$ 为半径：一个
约 $1.3$ 个像素的圆斑。

\textbf{7.} \qty{1}{''} $= \qty{4.9e-6}{rad}$，是极限的七倍；$D = 1.22\lambda/4.9
\times 10^{-6} = \qty{14}{cm}$ 即已达到视宁度。

\textbf{8.} \qty{2.7e-6}{rad}（\qty{0.55}{''}）；大气相位误差在较长波长中
所占比例较小：修正次数更少、速度更慢。

\textbf{9.} $0.15 < 0.55$：不能；在太空中于 \qty{550}{nm} 处为 \qty{0.14}{''}：恰好可以。

\textbf{10.} $6.7 \times 10^{-7} \times 3.8 \times 10^8 = \qty{250}{m}$；着陆器不能。

\textbf{11.} 四条衍射尖刺构成的十字形。

\textbf{12.} $(1000/7)^2 = 2 \times 10^4$：可见更暗的天体，而非更清晰的。

\textbf{13.} \qty{370}{nm} 与 \qty{240}{nm}：
分别对应每毫米 $2700$ 与 $4200$ 个线对。

\textbf{14.} 细菌可以；其 \qty{0.2}{\micro m} 结构处于极限，
需浸油，干式不行；病毒不能。

\textbf{15.} 在正常照明下 $d \ge \lambda/\mathrm{NA} = \qty{390}{nm}$，斜射光下为 \qty{200}{nm}
——两点的 $0.61\lambda/\mathrm{NA}$ 介于其间。

\textbf{16.} 油的折射率提高了 $n\sin\alpha$，
并抑制了会把陡峭光线困在盖玻片内的全反射。

\textbf{17.} \qty{110}{nm}；电子：$0.61 \times 4\,\text{pm}/0.01 = \qty{0.24}{nm}$。

\textbf{18.} 波片使未衍射的光延迟四分之一波长，
从而与经细胞相位结构衍射的光发生干涉；
相位图样变为强度图样。

\textbf{19.} \qty{4.3e-4}{rad}；$1.5 + 2 \times 43 = \qty{88}{mm}$。

\textbf{20.} $2.44\lambda f_1/D = \qty{8.7}{\micro m}$；针孔 \qty{13}{\micro m}。

\textbf{21.} $f_2 = 20f_1 = \qty{200}{mm}$；\qty{2.2e-5}{rad}；
行进一千米后为 $30 + 44 = \qty{74}{mm}$。

\textbf{22.} $\lambda/\qty{100}{\micro m} = \qty{5.3e-3}{rad}$：
在针孔处偏离光轴 \qty{53}{\micro m}——被挡住。

\textbf{23.} $1\%$，被针孔板吸收。

\textbf{24.} 直径为 $D$ 的光束在 $\sim D^2/\lambda$ 距离内扩大一倍：
\qty{1.5}{mm} 对应 \qty{4}{m}，\qty{30}{mm} 对应 \qty{1.7}{km}。

\textbf{25.} 眼睛 \qty{3}{mm}，望远镜 \qty{1}{m}，物镜（NA $1.4$），光束 \qty{30}{mm}：
在每种情形中，$\lambda/D$ 决定最小的角度、最细的细节，
以及最慢的扩展。
\end{solution}
